%% ma: L12-B08-0180
%% lop: 12
%% chuong: 2
%% bai: 8
%% dang: 12.8.53
%% loai: TN4
%% muc: VDC
%% dap_an: "B"
%% nguon: "CĐ2.4.pdf giáo viên gửi 27/10/2026 (M07, Lê Minh Khôi – Nguyễn Anh Khôi, Chương II Chủ đề 4 Các bài toán về tâm tỉ cự), Câu 30"
%% kien_thuc: "Bài 8 (khoảng cách hai điểm, biểu thức tọa độ), Bài 6 (bình phương vô hướng, tâm tỉ cự), Bài 5 (giá trị nhỏ nhất của hàm số bằng đạo hàm), Bài 8 (tích vô hướng: a + b + c - 8 = DM.(1;1;1) không quá DM.căn 3)"
%% trang_thai: cho-duyet
%% ly_do: "Dạng mới 12.8.53 đề xuất, chờ giáo viên duyệt"
\begin{ex}
Trong không gian $Oxyz$, cho $A(3;-2;-1)$, $B(1;4;6)$, $C(3;38;-16)$ và $D(2;5;1)$. Điểm $M(a;b;c)$ là điểm trong không gian sao cho biểu thức \[P=MD^4-\dfrac{16}3MD^3-4MA^2-5MB^2-MC^2\] đạt giá trị nhỏ nhất. Tìm giá trị lớn nhất của biểu thức $Q=a+b+c$.
\choice{$Q_{\max}=8-5\sqrt3$}{\True $Q_{\max}=8+5\sqrt3$}{$Q_{\max}=8$}{$Q_{\max}=13$}
\loigiai{$\dfrac{4A+5B+C}{10}=(2;5;1)=D$ nên $4MA^2+5MB^2+MC^2=10MD^2+4DA^2+5DB^2+DC^2$.

Đặt $t=MD\ge0$: $P=t^4-\dfrac{16}3t^3-10t^2+$ hằng số; $P'=4t(t+1)(t-5)$ nên $P$ nhỏ nhất khi $t=5$.

$M$ thuộc mặt cầu tâm $D$ bán kính $5$: $Q=a+b+c\le2+5+1+5\sqrt3=8+5\sqrt3$.}
\end{ex}
